\documentclass[12pt]{article}

\usepackage{enumitem}
\usepackage{amsmath}

\begin{document}

An Exploration Dynamical Systems Chaos Argyris C02

\hfill

pag13

\hfill

\textbf{2 Preliminaries}

\hfill

If therefore, those cultivators of the physical science from whom the
intelligent public deduce their conception of the physicist . . . are led
in pursuit of the arcanas of science to the study of the singularities
and instabilities, rather than the continuities and stabilities of things,
the promotion of natural knowledge may tend to remove that prejudice
in favour of determinism which seems to arise from assuming
that the physical science of the future is a mere magnified image of
that of the past.

James Clerk Maxwell, 1873

\hfill

This book is an attempt to convey concepts of methods evolved in the field of nonlinear
dynamics to budding physicists and engineers and to illustrate them using
simple examples. The basis for these new ideas on dynamics is the topological or geometrical
view of temporal processes which leads to a representation in phase space.

\hfill

In the following, we shall restrict ourselves mainly to those non-linear processes of
evolution, the dynamics of which can be described by ordinary differential equations
or difference equations. We consider systems with the following characteristics:

\begin{enumerate}[label=\roman*]
\item they are finite-dimensional, i.e. each state is determined by a point in a finitedimensional
phase space;

\item they are differentiable, i.e. the temporal change of state is modelled by differentiable
functions;

\item they are deterministic, i.e. the equations contain no stochastic terms. Moreover,
the initial state determines the solution uniquely at least locally. This means in
particular that two trajectories can never intersect.
\end{enumerate}

The differential-topological method offers the opportunity of using the structure of
these differential equations to investigate them geometrically in phase space. What
is novel in our present-day understanding is that we do not concentrate on a single
trajectory but consider the solutions as a whole and analyse their stability globally.
In the case of dissipative systems, we observe the ensemble of all the trajectories
and classify the long-term behaviour via the stability behaviour of individual
groups. This type of analysis is called qualitative dynamics or, in mathematical
terminology, the modern theory of dynamical systems. The logical consequence is
the investigation of the stability of states of equilibrium or periodic solutions, for
example, and, moreover, the study of bifurcations which occur for critical control
parameters where the dynamics abruptly undergoes structural changes. The decisive
factor here is the phenomenon of instability. This plays the central role in the
theory of dynamical systems. Instabilities are responsible for the creation of new
temporal patterns as a result of bifurcations and furthermore for the occurrence of
irregular motions. The consequence is that slight changes in the initial conditions
or the system parameters can lead to unpredictable magnification effects. This revolutionary
and at the same time startling discovery is fundamental and undermines
some supposedly irrefutable principles of science (Cvitanovi´c et al., 2008); this is
the subject of the following section.

\hfill

Prior to this, we would like to point out explicitly which systems are excluded in the
following considerations as a consequence of the above-mentioned three properties.
These are, in particular, partial differential equations, systems with time delay
as well as differential equations with an additional stochastic term which occur
for example in synergetics, which we address briefly in section 6.8. In addition, it
would go far beyond the scope of this book to discuss dynamical systems which
are only piece-wise continuous or piece-wise smooth, i.e. systems which contain
jump discontinuities in the function or in the derivative. This includes, for example,
mechanical systems with dry friction or systems where shocks or pulses occur,
systems with electronic switches, biological systems with thresholds or economic
systems. In practice, such systems occur often and can cause complex dynamical
behaviour as well as intricate bifurcation scenarios. Recently, these systems have
become the subject of intensive research, see (Zhusubaliyev and Mosekilde, 2003;
di Bernardo et al., 2008).

\hfill

\textbf{2.1 Causality – Determinism}

\hfill

Classical physics was based on the assumption that the future is determined by
the present and that, in this way, an accurate knowledge of the present can reveal
the future. This belief in a fundamental limitless predictability was expressed in
perhaps the most elegant and comprehensible way by Laplace in 1814 through his
fictional demon:

\hfill

“Une intelligence qui pour un instant donn´e, connaˆıtrait toutes les forces
dont la nature est anim´ee, et la situation respective des ˆetres qui la composent,
si d’ailleurs elle ´etait assez vaste pour soumettre ces donn´ees `a
l’analyse, embrasserait dans la mˆeme formule, les mouvements des plus
grands corps de l’univers et ceux du plus l´eger atome: rien ne serait incertain
pour elle, et l’avenir comme le pass´e, serait pr´esent `a ses yeux.”
(Laplace, 1812)

\hfill

Anyone expressing the certitude of natural laws so pointedly, who, indeed, interprets
causality so narrowly as to mean determinism, is asking to be contradicted,
particularly at a time when the consequences of the quantum theory and the development
of atomic physics affect not only philosophy but our thought processes.

\hfill

This confrontation makes it clear how generation-dependent the principle of causality
or the law of cause and effect is and how it changes in the course of developments
in thinking and language. We must conclude that the concept of causality and our
understanding of it is linked with the concepts and the meanings associated with
them at anyone time! Formulations such as causality, causal law, causal explanation
or causal principle are always confusing as long as their explanations are not
transparent. In this connection, it appears helpful to begin with the historical development
of these concepts.

\hfill

The fact that causality sets up a link between cause and effect is, historically speaking,
relatively recent. For Aristotle, the originator of many scientific disciplines, the
word causa had a much more general meaning than nowadays. He expressed his
thoughts in the two pithy concepts “xxxxx”. With direct reference
to Aristotle, scholasticism, the philosophy of the Middle Ages, taught that
there are are four types of cause: 1. causa materialis (= material cause); 2. causa
formalis (= formal cause); 3. causa finalis (= final cause); 4. causa efficiens (=
efficient cause). In addition, Aristotle differentiated between internal and external
causes, the first two mentioned above being internal, the second two external. The
one we roughly associate with the word “cause” nowadays is causa efficiens, the
initiating or efficient cause.

\hfill

With the emergence of the Renaissance and the birth of modern scientific thinking
– usually associated with names such as Nicolaus Copernicus, Galileo Galilei
and Johannes Kepler – a change in the concept of causa also came about. As man
turned away from metaphysics and towards physics, i.e. towards quantification and
metrology, the word causa was linked with the material incident which preceded
and somehow caused the phenomenon to be explained. At first, it was held that if
the cause had been discovered, the natural phenomenon could be explained. In this
connection, cause was synonymous with explanation. Gradually, however, the understanding
took hold that it was the natural laws, the interconnection of cause and
effect, which clearly determine and thus explain natural processes. The culmination
of this development was the epoch-making work of Isaac Newton. Immanuel Kant,
who, in his “Critique of Pure Reason”, equated science as such with the science of
Newton, found a formulation for causality which was still valid in the 19th century:
“Wenn wir erfahren, daß etwas geschieht, so setzen wir dabei jederzeit voraus, daß
etwas vorhergehe, woraus es nach einer Regel folgt.”

\hfill

Einstein praised Newton’s outstanding intellectual achievement as “perhaps the
greatest progress in thinking that a single individual ever had the privilege to
accomplish”. His most eminent achievement was his comprehensive mathematical
theory of mechanics which remained the basis of scientific thinking until well into
the 20th century. It is hardly surprising that his general laws of motion – obeyed by
all objects in the solar system from the apple falling from the tree to the planets –
or their universal applicability fed the expectation that processes in nature were in
principle clearly determined, assuming they are known as a whole or in part after
careful study. Newton’s universe was a grand mechanical system which functioned
in accordance with exact deterministic laws so that the motion of a system in the
future could be calculated in advance from the state of this system at any given
time. If one is convinced that nature fundamentally behaves in this way, the next
logical step is the one Laplace formulated for his fictional supernatural intelligence
mentioned above: all the future as well as the past of our universe is calculable from
the precise knowledge of the position and velocity of all atoms at any one instant.

\hfill

This concept of pre-calculability in all eternity is not consistent with 20th century
physics and particularly not with quantum mechanics. It is not that physics
is in fundamental opposition to man’s longing for predictablity but that the two,
quantum mechanics and determinism, are incompatible with one another in their
philosophical approaches. The atomistic ideas of Democritus and Leucippus handed
down from classical antiquity assume that regular, ordered processes on a global
level derive their morphology from the irregular, random behaviour on a local level.
Such considerations are indeed plausible, as is substantiated by numerous examples
from everyday life. A fisherman, for example, fighting against the wind and the
waves, only has to ascertain the rhythm of the waves in order to be able to react to
them; he does not have to know the motion of each individual drop of water. If we
explain the processes we can perceive with our senses by the interaction of many
individual processes on the local level, is it not an inevitable consequence to regard
the overall regularities of nature as statistical? Statistical laws may provide reliable
statements on the whole set of possible events such as, for example, on mean values
and possible deviations, yet there will never be a totally reliable prediction for a
particular event; this will always remain indeterminate.

\hfill

In spite of such conceptual difficulties, we are continually setting up statistical laws in
day-to-day life as a basis for our practical actions. Engineers, for example – whether
they are acting in the design of aircraft, buildings or machines – cannot base their
work on precise loads or material data. They have to rely as a matter of course on
mean, that is to say statistical characteristic values. Yet nevertheless, when our attention
is drawn to such “semi-exact” regularities, we feel uneasy and consider them
less trustworthy. We would prefer either precisely definable processes in nature or,
on the other hand, chaotic, totally irregular ones. Is there an explanation for such an
attitude? One generally uses statistical laws when the physical system in question is
only partially known. The simplest example of this is the game of heads or tails. Since
no one side of the coin has an advantage over the other, we have to come to terms with
the fact that – when playing a large number of games – we can only predict one of the
two results with 50\% certainty. Games of dice are similar, except that the number of
possible results has increased to six. The likelihood of predicting a chosen number is
thus reduced to 1/6. If the die is thrown often enough, then the number of throws
resulting in a 1 is approximately one-sixth of all the throws. This is, of course, only
true if we assume a perfect die and an identical throwing technique; only then can
all six possible results be considered equally probable.What we assume is that these
conditions are satisfied approximately.

\hfill

In modern times, it was Robert Boyle who took up the idea from classical antiquity
by not only describing the material behaviour on the macroscopic level of observation
qualitatively as the result of the statistical behaviour of the molecules in a gas but by
also quantifying the well-known relationship between pressure and volume. His idea
was that the pressure, a quantity measurable on the macroscopic level, was built up
by the numerous impacts of molecules on the wall of the vessel; this interpretation
was anmiraculous brainwave at the time. In a similar way, the laws of thermodynamics
could be conceived when it proved possible to state in a mathematically precise
form the fact that atoms move more violently in hot bodies than in cold ones.

\hfill

While Laplace raised the hopes that one day, the whole world could, in principle,
be calculable, in the second half of the last century, the idea spread that Newton’s
mechanics might be unrestrictedly valid, but that systems within the kinetic gas
theory could never be completely determinable due to the immense number of gas
molecules. It was mainly Josiah Willard Gibbs and Ludwig Boltzmann who put
the incomplete knowledge of these systems into mathematical language by using
statistical laws. Gibbs went one step further by introducing temperature as a physical
concept for the first time; this only makes sense when the knowledge of the
system is incomplete. If indeed the velocity and the position of all the molecules in
a gas were known, it would be completely superfluous or meaningless to talk of the
temperature of the gas. Hence, the concept of temperature can only be used meaningfully
when the system on the microscopic level of observation is incompletely
known and one nevertheless does not want to forgo a “qualitative” statement on
the macroscopic level. Using such a concept, one does not describe the behaviour
of a system by taking an increasing number of degrees of freedom into account –
in the ideal case, infinitely many – but by proceeding to new, essential and thus
considerably fewer degrees of freedom on a more general level of observation. The
motto is not “more precise, more detailed and infinitely many” but “more global,
fewer and nevertheless informative”. In this highly simplified statement, we share
the view of Max Born (1959) who says that absolute accuracy is not a physically
meaningful concept and can only be found in the conceptual world of mathematicians.
Felix Klein called for the application of “approximation mathematics” side by
side with the usual “precision mathematics”. Since his suggestion remained without
response at the time, the physicists at the turn of the preceding century solved their
problems in their own conceptual framework, using methods of probability and statistical
laws.What Max Born’s statement basically implied was that non-linear laws
and deterministic equations may supply unpredictable answers. At the same time,
he in effect pointed out that non-linear equations often react with unexpected sensitivity
to the slightest changes in the initial conditions and thus suddenly supply
unexpectedly differing answers. This was a revolutionary perception, first expressed
by Poincar´e at the turn of the preceding century but unnoted for a long time:

\hfill

“Une cause tr`es petite, qui nous ´echappe, d´etermine un effet consid´erable
que nous ne pouvons pas ne pas voir, et alors nous disons que cet effet est
dˆu au hasard. Si nous connaissions exactement les lois de la nature et la
situation de l’univers `a l’instant initial, nous pourrions pr´edire exactement
la situation de ce mˆeme univers `a un instant ult´erieur. Mais, lors mˆeme
que les lois naturelles n’auraient plus de secret pour nous, nous ne pourrons
connaˆıtre la situation initiale qu’approximativement. Si cela nous permet
de pr´evoir la situation ult´erieure avec la mˆeme approximation, c’est tout ce
qu’il nous faut, nous disons que le ph´enom`ene a ´et´e pr´evu, qu’il est r´egi par
des lois; mais il n’en est pas toujours ainsi, il peut arriver que de petites
diff´erences dans les conditions initiales en engendrent de tr`es grandes dans
les ph´enom`enes finaux; une petite erreur sur les premi`eres produirait une
erreur ´enorme sur les derniers. La pr´ediction devient impossible et nous
avons le ph´enom`ene fortuit.” (Poincar´e, 1908)


\hfill

In the first two decades of the 20th century, the attempt was initally made to explain
the motion of the atoms and molecules according to the basic precepts of
classical mechanics, in the spirit of Newton. The result, however, was an entanglement
of inextricable contradictions. For example, in accordance with classical
physics, a charged electron was supposed to orbit around the atomic nucleus and
constantly emit radiation until it collapsed into the nucleus due to loss of energy.
However, the electron paths devised in the model were unstable. It was obvious
to Niels Bohr that, on the basis of Planck’s theory, the paths of the electrons are
stationary, that the electron, as long as it sticks to its path, does not emit any radiation,
but that a change of path is accompanied by a loss of energy. Bohr solved
this contradiction not by changing his concept of the model, as could have been
expected, but rather by maintaining that classical physics was not applicable to the
description of the dynamical behaviour of atomic relations. The fact that loss of
energy which takes place in irregular bursts and in stages when an electron changes
its path should lead to the assumption that the radiation emittance of atoms is
a statistical phenomenon was acceptable. It is also possible to tolerate, at least
temporarily, a new type of mechanics, quantum mechanics, albeit at the expense of
determinacy, when the stability of the atoms can thus be mathematically ensured.
But the bold assertion that it is fundamentally impossible to know all the necessary
defining elements in order to achieve a complete determination of the processes –
even Albert Einstein could not and would not accept this as true, even for atomic
phenomena. For him, the quantum theory was only an instrument, temporarily
necessary; it had emerged due to our lack of knowledge of all the canonical variables
of the atomic process but could be suspended as soon as all these unknowns
had been clarified. His opinion of quantum mechanics culminated in the statement,
“God does not play dice”, to which Niels Bohr replied, “It would be presumptuous
of us human beings to prescribe to the Almighty how he is to take his decisions.”
What is it about the quantum theory that is so challengingly unfamiliar? It is the
fact that the concept of the trajectory has been banned from a mechanical theory
of the atom shell and replaced by irreducible probabilistic elements. This was
imperative in order to ensure discrete energy levels, whose existence had been revealed
by spectroscopy, in a mathematical equation for systems such as the atom.
It was Werner Heisenberg who took this radical step in 1925 in his decisive work
with the title “ ¨Uber quantentheoretische Umdeutungen kinematischer und mechanischer
Beziehungen” (Heisenberg, 1925). Here, the canonical variables position
and impulse become non-commutative quantities; Max Born and Pascual Jordan
recognised their matrix character. Although the fundamental features of this mathematical
theory had already been set up, Heisenberg described the problem of atomic
physics in a letter to Pauli in 1926 as completely unsolved. What caused him to
make this negative statement? What was still missing was the link to the physical
experiment (Heisenberg, 1969).

\hfill

The solution was provided, again by Heisenberg, by the indeterminacy relation in
1927 (Heisenberg, 1927). Position and velocity or momentum, the classical quantities
for the determination of a particle trajectory, can be stated individually with
arbitrary accuracy; simultaneously, however, this is impossible in the quantum philosophy, at least in the microworld of 
atomic physics. This is not an assumption
in quantum mechanics, but the consequence of the laws of the quantum theory.
It was thus clear that Newton’s mechanics, which assumes the exact knowledge of
both position and momentum in order to calculate a mechanical course of motion, is
not applicable to the atomic world. Although more than seventy years have passed
since Bohr, Heisenberg, Born and others came to the conclusion that the quantum
theory forces us to formulate the laws on the quantum level as statistical laws and
to bid farewell to determinism, it is difficult to incorporate this into our general
philosophy. For the atomic field, this call to abandon pure determinism may be
valid. But to speak of absolute randomness and to assert that the classical idea of
predictability is invalid in principle, although it was so successful in its search for
order, regularity and natural laws, merely due to spontaneous nuclear disintegration
for which there is neither cause nor explanation – this is something we still
cannot and do not want to bring ourselves to accept.

\hfill

In contrast to this, it should be remembered that classical physics tacitly complements
the principle of “identical” causality with the principle of “similar” causality.
Laplace’s principle, “identical causes have identical effects”, was extended by the
principle “similar causes have similar effects”. The reason is that it is a justifiable
expectation to demand exactly identical causes so as to guarantee the reproducibility
of physical processes, but that this is not feasible in practice. Any experimenter
knows that, although he is at pains to obtain the same results for repeated measurings
under the same test conditions, exactly identical repetitions of the test
conditions are basically impossible. The accuracy of measurements is limited, even
though it may be extremely high with today’s technical sophistication. When the
errors in the measurement results are of the same order as the inaccuracies of the
experimental set-up, they are not exactly the same, but similar. In spite of statistical
laws which we take as the basis of our experimenting, we still speak of the
reproducibility of physical behaviour. Due to this discrepancy between abstract
mathematical precision and unavoidable physical approximation, Max Born (1959)
demands the re-formulation of the question of determinability in mechanics and, in
its place, the differentiation between stable and unstable motions.

\hfill

Heisenberg and Born’s renunciation of the dogma of predictability may be valid
for the field of atomics, but appears strange when viewed on the macroscopic scale
which is directly accessible to our senses. We are, of course, aware that, in weather
forecasting, we have to rely on probability estimates although the motion of the
earth’s atmosphere follows exactly the same physical laws as the motion of the
planets. Nevertheless, the weather retains a considerable amount of randomness.
The reason for this is the exceptionally high number of unknowns necessary for a
detailed description of the dynamics involving correlations and feedback reactions
far beyond our knowledge. Until not so very long ago, there was little reason to
doubt that, in this case, at least in principle, accurate predictions would ultimately
be possible. In the spirit of Laplace’s demon, it was assumed that it was only necessary
to gather sufficient information about the system and to process it with the
necessary effort.

\hfill

This mechanical conception of the world which was first shaken by quantum mechanics
received a second blow as a result of an astonishing discovery, namely that
even simple non-linear systems may generate irregular behaviour. We have to come
to terms with the idea that such unpredictable behaviour is an intrinsic reality as it
does not disappear, even after more information has been gathered. Such seemingly
random behaviour which is generated in non-linear deterministic dynamics is called
“deterministic chaos”. The fact that deterministic laws without stochastic elements
cause chaos sounds as paradoxical as the concept of deterministic chaos itself. At
this point, we would like to stress that there are various manifestations of chaotic
behaviour between which we must carefully distinguish. In a vessel filled with gas,
the atoms fly around wildly and collide and only in this case, as Boltzmann already
established, does microscopic chaos prevail. In contrast to this, macroscopic or deterministic
chaos dominates when purely random oscillations occur although the
laws which describe the dynamics are deterministic.

\hfill

We have repeatedly mentioned the diametrically opposing conceptions of scientific
knowledge: on the one hand, the idea of the atomists who stress random collision
and, on the other hand, the mechanistic view of the world which is based on timeless
dynamic laws. Both conceptions fail when it is a question of explaining both spatial
and temporal structures, i.e. as occurring in undamped oscillations. Equilibrium
thermodynamics and conventional statistical physics do not provide the methods
to deal with such an oscillatory behaviour. A way out of this dilemma begins to
emerge.

\hfill

What is this hope based on? It is the study of the physics of non-equilibrium states
and on the other hand of the theory of dynamical systems. The physics of nonequilibrium
states deals with systems far from thermodynamic equilibrium where
chaos prevails on the microscopic level but is completely concealed on the macroscopic
level by well-organised patterns. In the second discipline, the theory of dynamical
systems, it is the instabilities that play the central role. At the point of
instability, the system must “choose” between various possibilities. A small, nonpredictable
fluctuation decides which course the dynamical process finally takes.
This occurrence of unpredictability forces us to re-think our understanding of deterministic
systems and the “long-term” prognoses we have won.

\hfill

Processes which are sustained by a continuous flow of energy and possibly matter
from outside and which are characterised by their ordered, self-organised, collective
behaviour are called evolutionary. A significant aspect of evolutionary processes is
their causal coherence, although they may be interspersed with random outbursts.
Anyone would reject the idea of regarding the random occurrences of a lottery as
an evolutionary development. Intuitively, we demand that the current state of a
system is fundamentally moulded, if not determined, by the preceding causes. This
does not mean, however, that this demand for causality totally excludes random
occurrences; evolutionary processes without mutation, symmetry breaking etc. are
unthinkable. Random occurrences reduce the chance of exact predictions, they cannot
be specified by laws, they weaken the deterministic net of cause and effect; yet
it would be wrong to draw the conclusion that the pattern of reality were chaotic.

\hfill

Thus, it is not only quantum physics, but also the chaos theory which throws considerable
doubt on Laplace’s demon of absolute determinism. “F¨ur manchen mag
dies eine Entt¨auschung sein. Aber vielleicht ist ja eine Welt sogar menschlicher, in
der nicht alles determiniert und nicht alles berechenbar ist, eine Welt, in der es –
dank der Quantenereignisse – Zufall, und damit auch Gl¨uck, gibt, in der – weil nicht
alle Probleme algorithmisch l¨osbar sind – Phantasie und Einfallsreichtum, Raten
und Probieren, Kreativit¨at und Originalit¨at noch gefragt sind und in der man, wie
die Chaos-Theorie zeigt, auch bei chaotischem Verhalten immer noch sinnvoll nach
einfachen Grundgesetzen suchen kann.” (Vollmer, 1988, p. 350)

\hfill

It is certainly a result of the historical development of the sciences that the unpredictable
(and thus also the chaotic state) was first explicitly suspected and then
formulated mathematically in the microworld of atomic physics. This does not mean
that these phenomena do not appear in our daily macroscopic life. Chaotic responses
of non-linear dynamical systems are almost a matter of course to us today. We draw
attention to the appearance of such phenomena in a pendulum with a large amplitude
and particularly so in the case of spatial oscillations. The well-known Duffing
equation is another example. A further historical example is the well-known anorganic
chemical reaction of Belousov-Zhabotinsky which offers a memorable display
of colour. Yet the most impressive and still the most mysterious is the onset and
development of turbulence in a fluid. Research into this problem – which is highly
important in technology and of decisive significance in meteorology and combustion
processes – has inspired many famous physicists since Osborne Reynolds to
look for a physical-mathematical solution. In spite of admirable contributions of
many scientists, even Heisenberg (Heisenberg, 1948) did not succeed in solving the
problem.

\hfill

For fully developed turbulence, however, today’s chaos theory is still not broadly
enough developed. In Chapter 9, we attempt to give an overview of the current
status of turbulence research and to show that further extensive physical considerations
and mathematical tools are necessary to model turbulent phenomena. Apart
from the difficulties involved in an adequate comprehension of a complex process,
there is a further obstacle that everyone breaking new ground must surmount, even
– or especially – in science. No one butWerner Heisenberg could have expressed this
so memorably: “Wenn wirkliches Neuland betreten wird, kann es aber vorkommen,
daß nicht nur neue Inhalte aufzunehmen sind, sondern daß sich die Struktur des
Denkens ¨andern muß, wenn man das Neue verstehen will. Dazu sind offenbar viele
nicht bereit oder nicht in der Lage”, see (Heisenberg, 1969), p. 102.

\hfill

\textbf{2.2 Dynamical Systems – Examples}

\hfill


Before we move on to a detailed description of the chaotic behaviour of systems
possessing deterministic equations of motion, we should like to lead up to this subject
by beginning with four typical examples.

\hfill

Our first example is a mechanical system with a single degree of freedom and no
loss of energy due to friction. Let a small mass be hung on a spring and allowed to
oscillate vertically (fig. 2.2.1). For small displacements, the motion of this undamped
spring pendulum can be expressed by the linear differential equation in the two
variables displacement x and its second derivative with respect to time t

\[
\ddot{x} + w^2_0x = 0
\]


where $w_0$ is the frequency of the oscillation.

Equation (2.2.1) has the general solution

\[
x = A \sin w_0 t + B \cos w_0t
\]

where the integration constants A and B are determined by the initial conditions.

Denoting the initial values of x and $\dot{x}$ at the instant t = 0 by $x_0$ and $\dot{x}_0$, we obtain
for the displacement x

\[
x = \frac{\dot{x}_0}{w_0} \sin w_0t + x_0 \cos w_0t
\]

and for the velocity $\dot{x}$

\[
\dot{x} = \dot{x}_0 \cos w_0 t - x_0 w_0 \sin w_0 t
\]

Eliminating the time t from eqs. (2.2.3) and (2.2.4), we find

\[
x^2 + \frac{\dot{x}^2}{w_0^2} = x_0^2 + \frac{\dot{x}_0^2}{w_0^2}
\]

The family of all ellipses in the $x,  \dot{x}$ -plane designated by eq. (2.2.5) is called a phase
portrait. On the right of fig. 2.2.1, the displacement x as a function of time for
a concrete initial condition is shown and on the left, the trajectories in the twodimensional
phase space for three different initial conditions. The state variables
or coordinates x and x˙ which span the phase space characterise the single-mass
system uniquely.

\hfill

In this idealised case where the energy is preserved, the mass returns periodically
to its initial position. Also, the states of maximum deflection and zero velocity are
repeated with a period of $T = 2 \pi/w_0$.

\hfill

In natura, it is practically impossible to realise the assumption “no loss of energy
due to friction”. Loss of energy, for example due to air resistance, continuously
decreases the deflection of the pendulum until at some stage, the mass remains
in a state of equilibrium. During the transient phase, the aforementioned elliptical
trajectories in the phase space then become spirals which all end in a point, the
state of equilibrium (fig. 2.2.2). This point which captures all spirals is called an
attractor, in this particular case a point attractor.

\hfill

When the loss of energy of the pendulum due to friction is of a viscous nature,
it can be reproduced by a damping term which depends to a first approximation
linearly on the velocity. The linear differential equation (2.2.1) for a conservative
system then takes the following standard form, again a linear system,

\[
\ddot{x} + 2 \zeta w_0 \dot{x} + w^2_0x = 0
\]


where the damping factor $\zeta > 0$ controls the fading of the periodic transient response.
In the case of damping $\zeta  > 1$, the system tends aperiodically towards the
state of equilibrium x = 0. For $0 < \zeta < 1$, the case of sub-critical damping, the
amplitudes also decrease; but the motion retains qualitatively the appearance of an
oscillatory process.

\hfill

The loss of energy in a system can be compensated by means of the continuous supply
of energy, e.g. a periodic external excitation (fig. 2.2.3). After a certain transient
phase, the trajectories in the two-dimensional phase space approach asymptotically
a closed curve, the so-called limit cycle, which is run through periodically. Beside the
point attractor, the limit cycle is the only possible attractor in the two-dimensional
phase space.

\hfill

For a harmonically excited, viscously damped oscillator, the equation of motion in
linearised form becomes

\[
\ddot{x} + 2 \zeta w_0 \dot{x} + w_0^2 x = w_0^2 x_0 \sin w_E t
\]


Here, the periodic external excitation is reproduced by a sinus term on the righthand
side. Linear equations of this type are integrable and lead to a family of
solutions, the individual curves of which are determined by the initial conditions.
Since the dependence of a specific solution is relatively insensitive with regard to
the initial conditions, small changes in the latter cause only small changes in the
solution. Thus, our physical system is characterised by a “similar” connection between
cause and effect.

\hfill

The situation in the case of non-linear equations of motion is completely different.
A simple dynamical system which leads to chaotic patterns of motion is the Duffing
equation (see section 10.5). In its modified form, the restoring force is approximated
by a polynomial of the third order, thus allowing, for example, the reproduction of
larger deflections in an externally excited beam under shear and normal force. In
this special case, the differential equation takes the form

\[
\ddot{x} + 2 \zeta w_0 \dot{x} - w_0^2 x + \delta w_0^2 x^3 = w_0^2 x_0 \sin w_E t
\]

where the cubic term $x^3$ possesses a stiffening effect. Due to its simple morphology,
the Duffing equation and its modifications led to numerous investigations on
the behaviour of dynamical systems both in mechanics and electrical engineering
(Ueda, 1980a; Ueda, 1980b; Moon and Holmes, 1979; Seydel, 1980). We stress that
the non-linearity – in x, but not in time t – is a necessary condition for chaotic
behaviour, though not sufficient in itself.

\hfill

If the substitution $x_1=x, x_2=\dot{x}, x3=t$ is carried out, the non-autonomous
differential eq. (2.2.8) with the second derivative with respect to time, the righthand
side of which depends explicitly on the time t, can be re-written as a system
of non-linear, so-called autonomous first-order differential equations

\[
\begin{aligned}
\dot{x}_1 &= x_2 \\
\dot{x}_2 &= -2 \zeta w_0 x_2 + w_0^2 x_1 - \delta w_0^2 x_1^3 + w_0^2 x_0 \sin (w_E x_3) \\
\dot{x}_3 &= 1
\end{aligned}
\]


This autonomous set of equations is equivalent to the non-autonomous eq. (2.2.8).
As a generalisation, we can state that an autonomous differential equation may be
put in the form

\[
\dot{\textbf{x}} = \textbf{F(x)}
\]

where the non-linear vector function \textbf{F(x)} does not depend explicitly on the time.
Thus, the introduction of additional variables transforms the non-autonomous
eq. (2.2.8) with second-order derivatives with respect to time into a system of
three first-order differential equations. Such a set of equations which is expressed
directly in terms of displacement and velocity considerably simplifies a qualitative
and quantitative discussion of the trajectories in phase space (the concept of the
phase space is elucidated in section 2.3).

\hfill

The phase portrait of eq. (2.2.9) for the variables $x_1, \dot{x}_1$ shows intersecting trajectories
corresponding to specific choices of control parameters and initial conditions
(fig. 2.2.4). Trajectories of the autonomous type of eq. (2.2.9) which do not intersect
in the extended phase space with the coordinates $x_1, \dot{x}_1$ and t do so in the projection
onto the ($x, \dot{x}$ )-plane, the customary phase space. It is easy to imagine that for a
non-periodic solution (i.e. an irregular or chaotic motion) for $t \to \infty$, the phase
portrait will be blackened to unrecognisability in a relatively short time. For this
reason, this representation must be ruled out as an indicator of chaos since it would
be futile to try to identify the corresponding trajectories which remain adjacent
over a long time but diverge exponentially from one another at some stage.

\hfill

The most widespread form of a display of the equations of motion consists of plotting
the deflection, for example, over time (fig. 2.2.5). However, to deduce a chaotic
behaviour by means of this method is not realistic and is doomed to failure because
the observed period of time is necessarily finite. Although the bizarre, non-periodic
course of the curve in fig. 2.2.5 strongly suggests an erratic course of motion, it
remains uncertain whether a periodic motion might not be established after all in
case one observed a longer time interval.

\hfill

It may require skill and effort to distinguish between regular and irregular behaviour
on the basis of the transient response of the variables of state of non-linear dynamical
models; it is, however, even more difficult to differentiate in an experimentally
established transient response between background noise and deterministic chaos.

\hfill

In order to assess dynamical behaviour, it is also possible to apply, apart from the
representation in phase space and the temporal change of individual variables of
state, the power spectrum method, well known to the engineer (see section 3.8).
The recorded time series, which are evaluated in this method by means of a Fourier
analysis, yield clearly defined peaks in the power spectrum diagram for periodic or
quasi-periodic motions; in contrast, continuous curves or curves with a high noise
level indicate chaotic or stochastic behaviour. This means that in the erratic case,
the measured variables can no longer be represented as a discrete superposition of
oscillations: from a critical value of this state variable in the Fourier-transformed
signal onwards, a very high noise component in the spectrum emerges. Since these
power spectra can easily be determined experimentally, however, they offer themselves
as a means of characterising the irregular behaviour of real systems. We can
illustrate this with an example taken from fluid dynamics (fig. 2.2.6). As a consequence,
simple spectral analysis does not suffice as a single tool for describing
erratic behaviour in chaotic physical systems. Other methods of measuring must
be found in order to describe the wide spectrum of regular and irregular motions
qualitatively.

\hfill

In recent years, a number of procedures have been proposed for noise reduction in
measured data sets (see (Kantz and Schreiber, 1997) and references cited therein).
Of particular interest is the method developed recently for analysing stochastic
data which is based on the theory of Markov processes and which allows the determination
of the model equations underlying the experimentally measured data
(Friedrich and Peinke, 1997; Friedrich et al., 2000). In contrast to other methods,
this is a parameter-free approach for reconstructing the model equations, i.e. a
priori, because no assumption needs to be made about the type of these equations.
Moreover, in many cases different types of noise can be separated from
the deterministic part (Siefert and Peinke, 2004; B$\ddot{o}$ttcher et al., 2006; Lehle, 2011;
Friedrich et al., 2011). As a result, one obtains stochastic differential equations
which contain both a deterministic and a stochastic part. However, in the following,
we restrict our considerations to dynamical systems which can be described solely
by deterministic equations.

\hfill

In the literature, a multitude of possible methods for characterising the occurrence
of chaotic motions is presented. Apart from the power spectrum and autocorrelation,
the two classic tools, we propose to concentrate on the following criteria:
Lyapunov exponents, dimensions and Kolmogorov-Sinai entropy. An extensive discussion
of these concepts can be found in the context of dissipative systems and
attractors in Chapter 5.

\hfill

The examples of mechanical systems mentioned here make it clear that the concept
of the attractor plays a central role in the description of the behaviour of damped
systems subject to deterministic equations of motion. We basically differentiate
between two types: regular attractors and strange attractors. On the one hand,
there are three classic types of motion: equilibrium, periodic motion and quasiperiodic
motion. All three states are associated with regular attractors since, in
the case of damping, the system tends towards one of these three states after the
transient phase. The point attractor corresponds to the state of equilibrium, the
limit cycle to periodic motion and the torus to quasi-periodic motion. On the other
hand, there is a class of deterministic, but erratic, i.e. chaotic, motions which are not
predictable if the initial conditions are subject to small fluctuations. Such long-term
behaviour ($t \to \infty$) is associated with the concept of the strange attractor.

\hfill

\textbf{2.3 Phase Space}

\hfill

In the case of evolutionary processes which can be analysed quantitatively, we
can record the modification of particular quantities as a function of time. These
variations must stem from specific causes. If the state is defined completely at a
given instant by n variables $x_1, … , x_n$, the evolution of the processmay be described
by a system of n ordinary differential equations

\[
\begin{matrix}
\frac{dx_1}{dt} = F_1(x_1, x_2, ... , x_n) \\
\frac{dx_2}{dt} = F_2(x_1, x_2, ... , x_n) \\
\vdots \\
\frac{dx_n}{dt} = F_n(x_1, x_2, ... , x_n)
\end{matrix}
\]


The n time-dependent variables represent physical quantities such as location, velocity,
temperature, pressure etc. Not only in physics, but also in other branches
of science such as biology, chemistry, economic sciences and other fields, many real
processes in time can be described by systems of ordinary differential equations. If
we formally introduce the column vectors

\[
\textbf{x} = \{ x_1 x_2 …  x_n \}
\]
\[
\textbf{F} = \{F_1 F_2 …  F_n \} 
\]

we can write the system eq. (2.3.1) in a simplified form (see also eq. (2.2.10))

\[
\frac{d \textbf{x}}{dt} = \textbf{F(x)}
\]

where n is the number of equations defining the whole system. The fact that the
system consists exclusively of first-order differential equations is not a restriction
since any system of ordinary differential equations of higher order can be transformed
into a system of first order by the introduction of additional variables.

\hfill

We wish to stress once again that the system eq. (2.3.1) or (2.3.3) is autonomous
since the right-hand side does not depend explicitly on the independent variable t.
This is not a restriction either, as each non-autonomous system of equations can be
transformed into an autonomous one by introducing an additional variable $x_{n+1} = t$
and the trivial relationship $\dot{x}_{ n+1} = 1$.

\hfill

It is very useful to represent the temporal evolution of a system in an abstract
space, the phase space. It is n-dimensional and is spanned by the state variables or
coordinates $x_1, x_2, …  , x_n$. In the phase space, the state of the system at a given
time is represented by a point. This point moves with time and its velocity is specified
by the vector \textbf{F}. Since the velocity \textbf{F} is known from eq. (2.3.3), the velocity
field can be represented immediately in the phase space (see fig. 2.3.1). The whole
set of directions specifies a direction field, reminiscent of the stream lines in a liquid.

\hfill

A point chosen arbitrarily at the instant $t = t_0$ describes a curve in phase space
which at each point is tangential to the vector field of velocities. The graph of the
motion in the phase space is called trajectory, phase line or orbit and the total
set of possible motions is denoted phase flow $\phi t$. The autonomous set of first-order
equations (2.3.1) allows us to display the velocity field without integration directly
in the phase space. We thus already obtain a first impression of the form of the
solution.

\hfill

Through each point of the phase space, exactly one trajectory runs. Physically
speaking, this means that if a state is known at a particular instant, both the
future and the past are determined by integration. This also means that trajectories
which represent a unique solution can never intersect. In the special case
of the mechanics of a particle where the dynamical state of a mass point is specified
in three-dimensional space by its position (three space coordinates) and by
its velocity (three velocity components), the phase space is six-dimensional. For a
single-degree-of-freedom oscillator, the phase space degenerates to the phase plane
with the components displacement and velocity (fig. 2.2.1).

\hfill

\textbf{2.4 First Integrals and Manifolds}

\hfill

In the distant past, when mathematical tradition demanded the solution of differential
equations by analytical integration, the somewhat strange concept of “first
integrals” was coined. What was then understood by an integral is nowadays called
solution.

\hfill

The definition of the first integral used here has been adopted from Arnold
(Arnold, 1980). Let \textbf{F} be a vector field, and the individual components $F_1, F_2 . . . F_n$
differentiable functions.

\hfill

\textbf{Definition:} a function \textbf{I(x)} is called the first integral of the differential equation

\[
\dot{\textbf{x}}  = \textbf{F(x)} 
\]

if its so-called Lie derivative $L_{\textbf{F}}$  (Olver, 1986) along the vector field \textbf{F} vanishes

\[
L_F \equiv F_1 \frac{\partial I}{\partial x_1} + ... + F_n \frac{\partial I}{\partial x_n} = \textbf{F}^t \frac{\partial 
I}{\partial \textbf{x}^t} = \frac{\partial I}{\partial \textbf{x}} \textbf{F}
\]

where the notation adopts the following conventions:

\[
\frac{\partial I}{\partial \textbf{x}} = \left [ \frac{\partial I}{\partial x_1} \frac{\partial I}{\partial x_2} 
...\frac{\partial I}{\partial x_n} \right ]
\]

denotes a row vector and, correspondingly,

\[
\left (  \frac{\partial I}{\partial \textbf{x}} \right )^t =  \frac{\partial I}{\partial \textbf{x}^t} =\left\{ \frac{\partial 
I}{\partial x_1} \frac{\partial I}{\partial x_2} ...\frac{\partial I}{\partial x_n} \right\}
\]

a column vector. Equation (2.4.2) implies the following characteristics of the first
integral: on the one hand, the function $I(x_1, x_2 . . . x_n)$ remains constant along each
given trajectory; on the other hand, each trajectory lies on a hypersurface in the
phase space. The hypersurface is defined by $I(\textbf{x}) = C$, where C is a constant (see
fig. 2.4.1 for a three-dimensional phase space). Each trajectory defined by the initial
condition thus lies entirely on a smooth hypersurface or generates it.

\hfill

The hypersurface defined by I(\textbf{x}) = C designates an (n - 1)-dimensional manifold (see also the discussion at the end 
of this section). If all possible values are assigned
to the parameter C, this leads to a one-parameter family of manifolds which fills
the phase space completely. If a first integral I is known, the equation I(x) = C
for a given C can be solved with respect to xi. If this xi is substituted into the
remaining (n - 1) equations of eq. (2.3.1), the whole set of equations is reduced
to (n - 1) equations. The more first integrals are known, the lower the number of
equations and the lower the dimension of the manifold generated by the trajectory.

\hfill

In mechanics, it is the conservation theorems which often supply first integrals; unfortunately,
however, no systematic rule is known which yields an easier derivation
of first integrals. The Hamilton function which represents a first integral for conservative
Hamilton systems is a stroke of luck. In addition, it can be shown that in
the case of conservative forces, the Hamilton function – if it is not time-dependent
– corresponds to the total energy, i.e. the sum of kinetic and potential energy (see
also section 4.1).

\hfill

It is necessary to associate many phenomena with geometrical models which cannot
be described in a simple form. In the case of dynamical systems, the geometrical
models form manifolds. Manifolds play an important role in the determination of
the global behaviour of trajectories. For example, non-chaotic attractors which specify
the asymptotic state of many dynamical systems lie on manifolds or mould them.

\hfill

An n-dimensional manifold M is a topological space with the property that each
point and its neighboorhood can be mapped by a one-to-one continuous bijective
mapping that has a continuous inverse function, a so-called homeomorphism, onto
an open subset of the n-dimensional Euclidean space E. Therefore, it is possible
to use the coordinates of E as local coordinates on M, see (Abraham et al., 1993;
Br¨ocker and J¨anich, 1990) and fig. 2.4.3. The underlying idea is to map globally
complex, but still sufficiently smooth structures onto elementary structures in order
to apply, for example, methods from differential geometry. It is well known
that the earth’s surface cannot be treated as a Euclidean space; measurements on
it must be performed in accordance with the laws of spherical trigonometry. Yet
parts of the earth can be mapped on charts of an atlas and, when approaching a
chart’s boundary, one has to switch to the next one.

\hfill

Geometries which may not admit an overall analytical representation are considered
locally in their tangential plane, i.e. the neighborhood of each point can be
described by local coordinates. By joining overlapping neighborhoods, so-called
charts, the entire manifold and the corresponding mapping can be described analytically.
We illustrate this approach using the limit cycle in fig. 2.4.2 as a simple
example of a manifold. The limit cycle of fig. 2.4.2 can be subdivided into an assembly
of overlapping subsets which can be mapped onto subintervals between 0 and $2|pi$a by a continuous one-to-one map. Thus, the 
$\phi$-coordinate can be used directly as a local coordinate on the manifold, independent of the phase space coordinates
which describe the limit cycle.

\hfill

A further example of a manifold in three dimensions is the two-dimensional torus
(fig. 2.4.5). For such a torus, each surface element can be mapped one-to-one onto
an element in the plane. The assembly of all the overlapping surface elements once
again forms a torus. In a corresponding manner, the surface elements on the associated
plane can be joined together. Their assembly yields the square in fig. 2.4.6.

\hfill

With scissors and glue and by rolling and bending, this square can be turned back
into the torus; in this way, it becomes clear that each single point of the torus is
defined by the coordinates $\phi_1$ and $\phi2$. A “helical line” on the torus corresponds
to a general straight line on the mapping plane. Figure 2.4.6 demonstrates both
coordinate lines $\phi_1$ = const and $\phi_2$ = const. Since both the limit cycle and the
torus are hypersurfaces in the phase space and a velocity vector \textbf{F} exists at each
point, they form differentiable manifolds.

\hfill

\textbf{2.5 Qualitative and Quantitative Approach}

\hfill

There are basically two contrary approaches to the study and understanding of
dynamical systems. In the first case, we concretise a particular problem as a dynamical
system and gather as much information as possible about its behaviour.
The logical consequence is a complicated set of equations, particularly since the
equations must be formulated as realistically as possible in order to incorporate all
the involved effects. In the second case, we are interested in the characteristics of
dynamical systems in general and not in entering into detail. Here, too, we must
differentiate between two cases:

\begin{enumerate}[label=\roman*]
\item A mathematical approximation in the classical sense which evolves and elucidates
this new qualitative analysis of differential equations and develops it on
the basis of assumptions and strict argumentation. Qualitative approaches are
based on geometrical or topological methods which the great mathematician
Henri Poincar´e (1854 – 1912) first applied to the investigation of the stability
of our solar system. Topology, the study of qualitative geometry, has become
an indispensible mathematical tool for describing the behaviour of dynamical
systems in all their complexity.

\item An approximation or simulation in the sense of experimental mathematics, an
approach made possible by the computer. In this case, the aim is to arrive at
generally valid statements on dynamical systems on the basis of simple nonlinear
dynamical systems, the behaviour of which is studied numerically. The
choice of representative examples in dynamical systems has to be guided by
intuition and experience in order to discover characteristic patterns of response
and thus obtain informative results for a broader class of dynamical systems.
\end{enumerate}

\begin{thebibliography}{9}
\bibitem{texbook}
John Argyris, Gunter Faust
 (2015) \emph{An Exploration of Dynamical
Systems and Chaos
}, Springer..
\end{thebibliography}


\end{document}



